\section{Related Work}

Coordinated multi-point transmission and network MIMO established the performance potential of multi-cell cooperation, especially when joint precoding mitigates inter-cell interference \cite{Gesbert10}. Subsequent work developed distributed weighted minimum mean-square error (WMMSE) formulations that decompose weighted sum-rate maximization into alternating updates for receivers, weights, and beamformers \cite{Shi11}. Related distributed beamforming methods also employ dual variables and interference prices to manage local power constraints and inter-cell externalities \cite{Dahrouj10}.

One of the earliest distributed formulations of cooperative base-station transmission employed block-diagonalization precoding together with dual decomposition for distributed power allocation under per-antenna constraints \cite{Jungnickel08}. That work by Hadisusanto et al. established an important decomposition principle for cooperative cellular systems. In contrast, the present paper removes the need for fixed block-diagonalization structures and develops a localized beamforming architecture whose performance provably approaches the globally coordinated benchmark as the cooperation radius increases.

Clustered cooperation and cell-free massive MIMO architectures also exploit local structure, but they are typically motivated by implementation or scalability considerations rather than an explicit locality-to-optimality scaling law \cite{Nayebi17}. Our positioning is therefore complementary: the paper combines distributed optimization, localized signaling, and an asymptotic theorem that quantifies how the residual optimality gap shrinks as the cooperation neighborhood expands.

This paper generalizes the distributed decomposition philosophy introduced by Hadisusanto et al. \cite{Jungnickel08} from fixed precoding plus power loading to fully adaptive localized joint beamforming and power allocation.
